\section{并发}
\subsection{并发}
Rust 中的线程安全体系依赖两个核心内置特型：@std::marker::Send@ 和 @std::marker::Sync@。
\begin{itemize}[itemsep=0.25em, parsep=0pt]
    \item Send：表示类型的值可以安全地移动到另一个线程。
    \item Sync：表示类型的不可变引用可以安全地在多个线程之间共享。
\end{itemize}

\begin{code}[title=fork and join]
    \begin{verbatim}
use std::{io, thread};

fn process_files_in_parallel(filenames: Vec<String>) -> io::Result<()> {
    const NTHREADS: usize = 8;
    let worklists = split_vec_into_chunks(filenames, NTHREADS);

    // fork
    let mut thread_handles = vec![];
    for worklist in worklists {
        thread_handles.push(
            thread::spawn(move || process_files(worklist))
        );
    }
    // join
    for handle in thread_handles {
        // std::thread::Result<std::io::Reault<()>>
        handle.join().unwrap()?;  
    }
    Ok(())
}
\end{verbatim}
\end{code}

如果需要在多个线程之间共享同一份数据，应使用 @Arc<T>@ 来实现安全的多线程共享，
每次在创建线程时，使用 @data.clone()@ 传递给此线程。

\begin{code}[title=rayon]
    \begin{verbatim}
use rayon::prelude::*;

// "Do two things in parallel"
let (v1, v2) = rayon::join(fn1, fn2);

// "Do N things in parallel"
giant_vector.par_iter().for_each(|value| {
    do_thing_with_value(value);
});
\end{verbatim}
\end{code}

\begin{code}[title=channel]
    \begin{verbatim}
fn start_file_reader_thread(documents: Vec<PathBuf>)
    -> (mpsc::Receiver<String>, thread::JoinHandle<io::Result<()>>) 
{
    let (sender, receiver) = mpsc::channel();

    let handle = thread::spawn(move || {
        ... 
    });
    (receiver, handle)
}
\end{verbatim}
\end{code}

当发送端产生值的速度快于接收端接收和处理值的速度时，通道中未处理的消息会不断积压，
为了解决这种背压问题，可以使用同步通道（@std::sync::mpsc::sync_channel@）。

\begin{code}[title=pipeline]
    \begin{verbatim}
fn run_pipeline(documents: Vec<PathBuf>, output_dir: PathBuf)
 -> io::Result<()> {
    // Start all 5 stages of the processing pipeline
    let (texts, h1) = start_file_reader_thread(documents);
    let (pints, h2) = start_file_indexing_thread(texts);
    let (gallons, h3) = start_in_memory_merge_thread(pints);
    let (files, h4) = start_index_writer_thread(gallons, &output_dir);
    let result = merge_index_files(files, &output_dir);

    // Wait for all threads to finish
    let r1 = h1.join().unwrap();
    h2.join().unwrap();
    h3.join().unwrap();
    let r4 = h4.join().unwrap();

    // Return the first error encountered (if any)
    // Note: h2 and h3 won't fail
    r1?;
    r4?;
    result
}
\end{verbatim}
\end{code}

